\section{Results}
\label{sec:results}

\subsection{Closed-Set Retrieval}
\label{sec:results:retrieval}

We first report closed-set retrieval on the four held-out external datasets.
All models use identical gallery/query roles and the same five split seeds;
therefore, differences reflect the learned embedding rather than a change in
evaluation membership. Table~\ref{tab:closed_set_results} will report
per-dataset Rank-1, Rank-5, Rank-10, CMC@20, CMC@50, and mAP, followed by the
dataset-macro average. We will also report identity-balanced Rank-1 to
separate performance on frequent identities from performance on identities
with fewer observations.

\begin{table*}[t]
\centering
\caption{Closed-set retrieval on held-out external datasets. Values will be
filled from the frozen five-seed evaluation. Best values are shown in bold.}
\label{tab:closed_set_results}
\begin{tabular}{l l r r r r}
\toprule
\textbf{Dataset} & \textbf{Model} & \textbf{Rank-1} &
\textbf{Rank-5} & \textbf{mAP} & \textbf{ID-bal. Rank-1} \\
\midrule
\multicolumn{6}{c}{\emph{Results pending final evaluation}} \\
\bottomrule
\end{tabular}
\end{table*}

The aggregate comparison will be complemented by per-dataset analysis. In
particular, we will distinguish gains on datasets with repeated, relatively
stable views from gains under larger viewpoint, illumination, background, and
species shifts. Paired confidence intervals are computed on the same queries
and identity blocks, rather than inferred from differences between independent
aggregate means.

\subsection{Open-Set and Open-World Performance}
\label{sec:results:openworld}

The AnimalCLEF-style results come from a separate open-set discovery
protocol, not from the closed-set retrieval benchmark. The protocol first
creates observation-aware gallery, known-query, and novel-query roles. It
then computes WildFusion query--query and query--gallery similarities,
selects per-animal clustering configurations from development similarities,
and applies the frozen configurations with Chinese Whispers to held-out
matrices. The table below will report discovery metrics for each animal
setting, including gallery-based and gallery-free cases where applicable.

\begin{table*}[t]
\centering
\caption{AnimalCLEF-style open-set discovery evaluation. Values will be
filled after the final runs.}
\label{tab:open_world_results}
\begin{tabular}{l l r r r r r}
\toprule
\textbf{Animal setting} & \textbf{Protocol} & \textbf{ARI} &
\textbf{AMI} & \textbf{V-measure} & \textbf{B$^3$ F1} &
\textbf{BAUS} \\
\midrule
Per-animal & Frozen clustering protocol & -- & -- & -- & -- & -- \\
\bottomrule
\end{tabular}
\end{table*}

The frozen cluster assignments are compared with held-out labels only to
compute discovery metrics. Clustering configurations are not refit on the
evaluation matrices. This separation prevents the AnimalCLEF-style results
from becoming oracle threshold diagnostics.

\subsection{Ablation Studies}
\label{sec:results:ablations}

We will isolate the contribution of the category-level term, inverse-frequency
category weighting, conditional within-category identity supervision, and
identity sub-centers. Additional controlled analyses will measure the effect
of observation-aware partitioning, animal-focused preprocessing, gallery
aggregation, and development-only threshold calibration. Each ablation keeps
the split manifests and evaluation seeds fixed and changes one component at a
time.

\subsection{Qualitative Analysis and Failure Cases}
\label{sec:results:qualitative}

We will visualize ranked retrievals and open-world assignments for
representative successes and failures. The analysis will cover viewpoint and
occlusion changes, age or appearance changes, illumination and background
variation, visually similar individuals, segmentation failures, and
cross-dataset or cross-species domain shift. For online evaluation, examples
will also show when an early false discovery affects later reassociation.
